\section{GDPval: comparar directamente el output con expertos de la industria}

\subsection{Pasar de "logró completarlo" a "¿la entrega es suficientemente buena?"}

\videofigure{figures/economic-value.jpg}{GDPval no solo pregunta si el agent completó la tarea, sino que hace que expertos comparen la IA con los productos de trabajo reales de profesionales del mismo ramo.}{00:28:06--00:28:38}

\videofigure{figures/gdpval-sectors.jpg}{Las tareas cubren nueve sectores de alto PIB y 44 ocupaciones, en un intento de representar el trabajo de conocimiento común en la economía.}{00:28:38--00:30:02}

\videofigure{figures/gdpval-task-examples.jpg}{Los productos de las tareas incluyen CAD, análisis de inversión, evaluación de imágenes de enfermería, diapositivas, tablas y reportes operativos.}{00:30:02--00:34:16}

A diferencia de los benchmarks de código, los outputs de GDPval por lo general no tienen una única respuesta correcta. La evaluación usa pairwise preference de expertos de la industria: se compara el producto de la IA con el producto del experto humano sin saber cuál es cuál, y se determina si gana, pierde o hay empate.

\begin{importantbox}{GDPval mide la calidad del producto, no el proceso de pensamiento oculto}
Los expertos ven la entrega final, así que la métrica se acerca más a la experiencia del cliente; pero no puede indicar directamente si el proceso del agent es confiable, si citó hechos incorrectos o si fallaría con otra entrada.
\end{importantbox}

\subsection{Construcción de los datos y características de las tareas}

\videofigure{figures/gdpval-construction.jpg}{GDPval primero selecciona los sectores según el PIB, y luego elige las ocupaciones y las actividades de trabajo típicas, con una construcción representativa de arriba hacia abajo.}{00:34:16--00:35:08}

\videofigure{figures/gdpval-task-characteristics.jpg}{El tiempo promedio de finalización humana de una tarea es de aproximadamente siete horas; las tareas son multimodales, dependen de archivos de referencia y, en su mayoría, no tienen una única respuesta correcta.}{00:35:08--00:36:08}

Este tipo de tareas se acerca más al trabajo de oficina que las tareas de software de METR, pero también es más difícil de calificar automáticamente. Una evaluación de alta calidad requiere expertos de dominio, un rubric explícito, pruebas ciegas y verificaciones de consistencia, con un costo mucho mayor que ejecutar pruebas unitarias.

\begin{warningbox}{La evaluación de preferencia se ve afectada por el estilo y la presentación}
Un reporte con buena tipografía y lenguaje fluido puede ocultar errores factuales; un producto humano profesional pero sencillo también puede subestimarse. Se debe registrar por separado accuracy, instruction following, formatting y severity, en lugar de mirar solo la preferencia general.
\end{warningbox}

\subsection{El avance de los modelos no sigue una sola línea}

\videofigure{figures/gdpval-improvement.jpg}{El pairwise win rate en GDPval mejora de manera significativa entre generaciones de modelos, pero de forma escalonada, no como un crecimiento lineal suave.}{00:36:08--00:38:18}

\videofigure{figures/model-specific-strengths.jpg}{Distintos modelos muestran fortalezas claramente diferenciadas: unos destacan en estética y documentos, otros son mejores en cálculo, instruction following o herramientas complejas.}{00:38:18--00:39:08}

El ranking de modelos depende de la combinación de tareas. Un modelo puede dominar en slides, PDF y disposición visual, mientras que otro puede ser más fuerte en cálculo de tablas, precisión textual o herramientas de código. Por eso, el despliegue empresarial debería enrutar por tarea, en lugar de suponer que un solo modelo domina todas las ocupaciones.

\begin{knowledgebox}{El model routing es un producto directo de la evaluación}
Si el benchmark conserva etiquetas de grano fino sobre ocupación, modalidad y tipo de error, se puede entrenar un router: elegir el modelo, la cadena de herramientas y la intensidad de revisión según la tarea, en lugar de usarlo solo para publicar una tabla de posiciones general.
\end{knowledgebox}

\subsection{Severidad de las fallas y costo económico}

\videofigure{figures/gdpval-failure-modes.jpg}{Los problemas comunes incluyen errores de instruction following, formato y precisión; la composición de errores difiere entre modelos.}{00:39:08--00:40:08}

\videofigure{figures/failure-severity.jpg}{Muchas fallas son del tipo "aceptable pero no lo suficientemente bueno"; la proporción de errores verdaderamente catastróficos es baja pero no se puede ignorar.}{00:40:08--00:41:08}

\videofigure{figures/economic-speed-cost.jpg}{Al considerar "intentar varias veces y luego dejar que un humano lo corrija", el modelo puede aportar beneficios de velocidad y costo, pero el beneficio depende del costo de reintento y de revisión.}{00:41:08--00:42:02}

Se puede expresar con un modelo simplificado de costo esperado:

$$
C(n)=nC_{\mathrm{AI}}+
\bigl(1-P_n(\text{acceptable})\bigr)C_{\mathrm{repair}}+C_{\mathrm{review}}.
$$

\begin{itemize}
    \item $n$: número de generaciones o reintentos;
    \item $C_{\mathrm{AI}}$: costo de cada ejecución del modelo;
    \item $P_n$: probabilidad de obtener un resultado aceptable a partir de $n$ intentos;
    \item $C_{\mathrm{repair}}$: costo de que un experto repare un producto fallido;
    \item $C_{\mathrm{review}}$: costo de revisión necesario sin importar si hubo éxito o no.
\end{itemize}

\begin{importantbox}{El beneficio económico proviene del flujo de trabajo "IA + humano"}
Si la revisión y la reparación son baratas, puede haber valor incluso con una tasa de éxito baja en un solo intento; si los errores son difíciles de detectar o tienen consecuencias graves, el costo de la revisión humana anula la ventaja de velocidad del modelo.
\end{importantbox}

\videofigure{figures/performance-variation.jpg}{El desempeño en GDPval varía según el sector, la duración de la tarea y la modalidad; las tareas cortas y algunos sectores se acercan más al nivel experto.}{00:42:02--00:45:48}

\subsection{Resumen de la sección}

GDPval hace avanzar la evaluación de "responder correctamente una pregunta" a "entregar un producto profesional comparable". Revela las fortalezas de los modelos, la severidad de los errores y el costo de la colaboración humano-máquina, pero sigue estando influido por el contexto y las preferencias de los expertos.

